\chapter{Groups as Symmetries}

\begin{center}\itshape
A cube can be put back in its box in $24$ ways, and shuffling four cards can be done in $24$ ways.\\
These are the same $24$, in a sense this chapter makes precise.
\end{center}

In 1974 Ern\H{o} Rubik, a Hungarian lecturer in design, built a wooden cube whose faces could twist, to show his students how parts can move
independently. Its plastic descendant has $43\,252\,003\,274\,489\,856\,000$ positions, and in 2010 a team running a computation on donated
machines proved that every one of them can be solved in at most $20$ moves. Nobody could check those numbers by twisting. Both come from the
same source: the moves of the cube can be composed and undone, and that makes them a \emph{group}. This chapter is about what that word buys.

A symmetry of an object is a transformation that leaves it looking the same. Two symmetries can be performed one after the other, and the result is again a symmetry; every symmetry can be undone. These two facts, and nothing else, are what a \emph{group} records. In this chapter we define groups, meet the examples that will run through the whole volume, and prove the first theorem of the subject: the size of a subgroup divides the size of the group.

\section{Groups}

\begin{definition}\label{def:group}
A \emph{group} is a set $G$ with a map $G\times G\to G$, $(a,b)\mapsto ab$, such that
\begin{enumerate}[label=(\roman*)]
\item (associativity) $(ab)c=a(bc)$ for all $a,b,c\in G$;
\item (identity) there is $e\in G$ with $ea=ae=a$ for all $a\in G$;
\item (inverses) every $a\in G$ has some $b\in G$ with $ab=ba=e$.
\end{enumerate}
The group is \emph{abelian} if $ab=ba$ for all $a,b$. Its \emph{order} $|G|$ is its number of elements.
\end{definition}

For abelian groups the operation is often written $a+b$, the identity $0$ and the inverse $-a$.

\begin{proposition}[first rules]\label{prop:rules}
In a group $G$:
\begin{enumerate}[label=(\alph*)]
\item the identity is unique, and each $a$ has a unique inverse, written $a^{-1}$;
\item (cancellation) $ab=ac$ implies $b=c$, and $ba=ca$ implies $b=c$;
\item $(ab)^{-1}=b^{-1}a^{-1}$ and $(a^{-1})^{-1}=a$;
\item for $a,b\in G$ the equation $ax=b$ has exactly one solution, $x=a^{-1}b$. So in the multiplication table of a finite group, every row (and every column) contains each element exactly once.
\end{enumerate}
\end{proposition}
\begin{proof}
(a) If $e,e'$ are identities, $e=ee'=e'$. If $b,b'$ are inverses of $a$, then $b=be=b(ab')=(ba)b'=eb'=b'$.
(b) Multiply $ab=ac$ on the left by $a^{-1}$ and use associativity; the other side is the same on the right.
(c) $(ab)(b^{-1}a^{-1})=a(bb^{-1})a^{-1}=aa^{-1}=e$, and similarly in the other order, so $b^{-1}a^{-1}$ is the inverse of $ab$ by (a). The element $a$ is an inverse of $a^{-1}$.
(d) $x=a^{-1}b$ is a solution, and it is the only one by (b). The row of $a$ lists the products $ax$; each $b$ appears for exactly one $x$.
\end{proof}

Powers are defined by $a^0=e$, $a^{n+1}=a^na$ and $a^{-n}=(a^{-1})^n$; by induction $a^ma^n=a^{m+n}$ and $(a^m)^n=a^{mn}$ for all $m,n\in\Z$.

\begin{example}[the examples of this volume]\label{ex:examples}
\begin{enumerate}[label=(\arabic*)]
\item $(\Z,+)$, $(\Q,+)$, $(\R,+)$, $(\C,+)$, and $(\Q^*,\times)$, $(\R^*,\times)$, $(\C^*,\times)$: abelian.
\item $\Z/n\Z$ under addition, and $(\Z/n\Z)^*$, the invertible classes, under multiplication (Volume~1, Chapter~3). The group of $n$-th roots of unity $\mathbb U_n=\{z\in\C:z^n=1\}$ (Volume~0, Chapter~2).
\item The \emph{symmetric group} $\mathfrak S(X)$ of all bijections of a set $X$, under composition: $(\sigma\tau)(x)=\sigma(\tau(x))$. For $X=\{1,\dots,n\}$ it is written $\mathfrak S_n$, and $|\mathfrak S_n|=n!$. It is not abelian for $n\ge3$: $(1\,2)(2\,3)$ sends $3\mapsto1$, while $(2\,3)(1\,2)$ sends $3\mapsto2$.
\item The \emph{general linear group} $GL_n(K)$ of invertible $n\times n$ matrices over a field $K$, and its subgroup $SL_n(K)$ of matrices of determinant $1$ (Volume~0, Chapter~7).
\item The \emph{dihedral group} $D_n$ ($n\ge3$): the isometries of the plane that map the set of vertices of a regular $n$-gon to itself.
\end{enumerate}
\end{example}

In (3) the product $\sigma\tau$ means ``first $\tau$, then $\sigma$'', as for composition of functions. Associativity holds because composition of maps is associative, and that is the deep reason why all groups of symmetries are groups.

\begin{proposition}[the dihedral group]\label{prop:dihedral}
Let $P$ be a regular $n$-gon centred at $0$, with vertices $v_k=e^{2ik\pi/n}$ in $\C=\R^2$. Let $r$ be the rotation $z\mapsto e^{2i\pi/n}z$ and $s$ the reflection $z\mapsto\bar z$. Then
\[
D_n=\{e,r,r^2,\dots,r^{n-1},\ s,sr,sr^2,\dots,sr^{n-1}\},
\]
these $2n$ elements are distinct, $r^n=s^2=e$, and $srs=r^{-1}$.
\end{proposition}
\begin{proof}
First, an isometry $g$ of the plane with $g(0)=0$ is linear. It preserves norms ($|g(x)|=|g(x)-g(0)|=|x|$), hence scalar products, by $\langle x,y\rangle=\frac12\left(|x|^2+|y|^2-|x-y|^2\right)$. So $g(e_1),g(e_2)$ is an orthonormal basis, and the coordinates of $g(x)$ in it are $\langle g(x),g(e_i)\rangle=\langle x,e_i\rangle$: $g(x)=x_1g(e_1)+x_2g(e_2)$. Consequently every isometry $f$ is affine ($f-f(0)$ is linear), and affine maps preserve averages of points.
Now let $f\in D_n$. It permutes the vertices, so it fixes their average, which is $0$ (the sum of the $n$-th roots of unity vanishes). So $f$ is linear, and it is determined by the images of $v_0=1$ and $v_1$, which form a basis of $\R^2$. Now $f(v_0)$ is some vertex $v_k$, and $f(v_1)$ is a vertex at the same distance from $v_k$ as $v_1$ from $v_0$, so $f(v_1)=v_{k+1}$ or $v_{k-1}$. That gives at most $2n$ possibilities. The $2n$ listed maps realise them: $r^k$ sends $(v_0,v_1)$ to $(v_k,v_{k+1})$, and $sr^{k}$ sends them to $(v_{-k},v_{-k-1})$, since $sr^k(z)=e^{-2ik\pi/n}\bar z$. They are distinct because they act differently on $(v_0,v_1)$. Finally $r^n=e$, $s^2=e$, and $srs(z)=\overline{e^{2i\pi/n}\bar z}=e^{-2i\pi/n}z=r^{-1}(z)$.
\end{proof}

The relation $srs=r^{-1}$ lets us compute everything: $sr^k=r^{-k}s$, so any product of $r$'s and $s$'s can be pushed into the form $r^ks^\varepsilon$. The multiplication table of $D_4$ is drawn in Figure~\ref{fig:symmetries}.

\begin{think}
Before reading on, guess how many rotations of space map a cube to itself (the identity included). Is it $6$, $8$, $24$ or $48$? A first idea: a rotation can put a chosen face on any of the $6$ faces. Then what freedom is left?
\end{think}

\begin{proposition}[rotations of the cube]\label{prop:cube}
The rotations of $\R^3$ that map the cube $[-1,1]^3$ to itself form a group of order $24$. Apart from the identity, they are: rotations by $\pm90^\circ$ and $180^\circ$ about the $3$ axes through the centres of opposite faces ($9$ rotations), by $\pm120^\circ$ about the $4$ long diagonals ($8$), and by $180^\circ$ about the $6$ axes through the midpoints of opposite edges ($6$).
\end{proposition}
\begin{proof}
A rotation $f$ is a linear map with orthogonal matrix of determinant $1$. If it preserves the cube, it permutes the $8$ vertices $(\pm1,\pm1,\pm1)$ (the points of the cube farthest from $0$). Since $e_1=\frac12\left((1,1,1)+(1,-1,-1)\right)$, the vector $f(e_1)$ is the midpoint of two vertices, so its coordinates lie in $\{-1,0,1\}$; and $|f(e_1)|=1$ since $f$ is orthogonal. So $f(e_1)=\pm e_i$ for some $i$, and likewise for $e_2,e_3$; as the $f(e_j)$ are orthogonal, the $i$'s are distinct. So the matrix of $f$ is a \emph{signed permutation matrix}: each row and each column has one non-zero entry, equal to $\pm1$. Conversely each of the $3!\cdot2^3=48$ signed permutation matrices preserves the cube, and exactly half of them have determinant $1$ (changing the sign of the first column is a bijection between determinant $1$ and determinant $-1$). So there are exactly $24$ rotations of the cube. The rotations listed in the statement do preserve the cube and are distinct, and $1+9+8+6=24$, so the list is complete.
\end{proof}

\begin{figure}[ht]
\centering
\fig{ch01-symmetries}
\caption{Left: the multiplication table of $D_4$, the $8$ symmetries of a square ($r$ is the quarter turn, $s$ a reflection). Each row and each column contains every element once (Proposition~\ref{prop:rules}(d)). The four blocks show that rotation times rotation and reflection times reflection are rotations, while the mixed products are reflections. Right: the three kinds of rotation axes of the cube, with the number of non-trivial rotations of each kind.}
\label{fig:symmetries}
\end{figure}

\section{Subgroups and the order of an element}

\begin{definition}
A \emph{subgroup} of $G$ is a subset $H$ that contains $e$ and is stable under products and inverses. It is then a group for the same operation, and we write $H\le G$.
\end{definition}

\begin{proposition}\label{prop:subgroup}
\begin{enumerate}[label=(\alph*)]
\item A non-empty $H\subseteq G$ is a subgroup if and only if $ab^{-1}\in H$ for all $a,b\in H$.
\item An intersection of subgroups is a subgroup. For $S\subseteq G$, the intersection $\langle S\rangle$ of all subgroups containing $S$ is the set of finite products $s_1^{\pm1}\cdots s_k^{\pm1}$ with $s_i\in S$ (the empty product being $e$). It is the \emph{subgroup generated by $S$}.
\end{enumerate}
\end{proposition}
\begin{proof}
(a) A subgroup has the property. Conversely, take $a\in H$: then $e=aa^{-1}\in H$, $b^{-1}=eb^{-1}\in H$ for $b\in H$, and $ab=a(b^{-1})^{-1}\in H$.
(b) The first claim is clear from the definition. The set $W$ of such products contains $S$, contains $e$, and is stable under products (concatenate) and inverses (by Proposition~\ref{prop:rules}(c), $(s_1^{\pm1}\cdots s_k^{\pm1})^{-1}=s_k^{\mp1}\cdots s_1^{\mp1}$). So $W$ is a subgroup containing $S$, and every subgroup containing $S$ contains $W$. Hence $W=\langle S\rangle$.
\end{proof}

For instance $D_n=\langle r,s\rangle$, the rotations form the subgroup $\langle r\rangle$ of order $n$, and $SL_n(K)\le GL_n(K)$.

\begin{proposition}[order of an element]\label{prop:order}
Let $g\in G$. Either all the powers $g^k$, $k\in\Z$, are distinct, or there is a least $n\ge1$ with $g^n=e$, called the \emph{order} $o(g)$ of $g$. In the second case:
\begin{enumerate}[label=(\alph*)]
\item $g^k=e$ if and only if $n\mid k$, and $g^k=g^l$ if and only if $k\equiv l\pmod n$;
\item $\langle g\rangle=\{e,g,\dots,g^{n-1}\}$ has exactly $n$ elements;
\item $o(g^k)=\dfrac n{\gcd(n,k)}$.
\end{enumerate}
\end{proposition}
\begin{proof}
If two powers coincide, $g^k=g^l$ with $k>l$, then $g^{k-l}=e$ with $k-l\ge1$, and a least such exponent exists.
(a) Divide: $k=qn+r$ with $0\le r<n$. Then $g^k=(g^n)^qg^r=g^r$, which is $e$ if and only if $r=0$ by minimality of $n$. Apply this to $k-l$ for the second claim.
(b) By (a), the powers $g^k$ are the $g^r$ with $0\le r<n$, and these are distinct.
(c) Let $d=\gcd(n,k)$. By (a), $(g^k)^m=e$ iff $n\mid km$ iff $\frac nd\mid\frac kdm$ iff $\frac nd\mid m$ (Gauss' lemma, Volume~1, since $\frac nd$ and $\frac kd$ are coprime). The least such $m\ge1$ is $\frac nd$.
\end{proof}

In $D_n$, $o(r)=n$ and every $sr^k$ has order $2$ (by $srs=r^{-1}$, $(sr^k)^2=sr^ksr^k=r^{-k}r^k=e$). In $\mathfrak S_n$, a $k$-cycle has order $k$.

\section{Lagrange's theorem}

\begin{think}
In $\mathfrak S_3$, the subgroup $H=\langle(1\,2)\rangle=\{e,(1\,2)\}$ has $2$ elements. Multiply every element of $H$ on the left by each of the $6$ permutations of $\mathfrak S_3$. How many different sets $gH$ do you get, and how big is each? What does that suggest about $|H|$ and $|\mathfrak S_3|$?
\end{think}

\begin{definition}
Let $H\le G$. The \emph{left cosets} of $H$ are the sets $gH=\{gh:h\in H\}$, $g\in G$. Their number (possibly infinite) is the \emph{index} $[G:H]$.
\end{definition}

\begin{theorem}[Lagrange]\label{thm:lagrange}
Let $G$ be a finite group and $H\le G$. The left cosets of $H$ partition $G$, and each has $|H|$ elements. Hence
\[
\begin{gathered}|G|=[G:H]\cdot|H|,\qq\text{and in particular }|H|\text{ divides }|G|.\end{gathered}
\]
\end{theorem}
\begin{proof}
Every $g$ lies in its own coset $gH$, since $e\in H$. If two cosets meet, $g_1h_1=g_2h_2$, then $g_1=g_2h_2h_1^{-1}\in g_2H$, so $g_1H\subseteq g_2HH=g_2H$, and symmetrically $g_2H\subseteq g_1H$: two cosets are equal or disjoint. Finally $h\mapsto gh$ is a bijection $H\to gH$ (its inverse is $x\mapsto g^{-1}x$), so $|gH|=|H|$. Counting $G$ coset by coset gives $|G|=[G:H]\,|H|$.
\end{proof}

\begin{corollary}\label{cor:lagrange}
Let $G$ be finite of order $n$.
\begin{enumerate}[label=(\alph*)]
\item The order of every element divides $n$, and $g^n=e$ for all $g\in G$.
\item (Euler--Fermat, again) For $a$ coprime to $m$, $a^{\varphi(m)}\equiv1\pmod m$.
\item If $n$ is prime, $G$ is generated by any of its elements $g\ne e$.
\end{enumerate}
\end{corollary}
\begin{proof}
(a) $o(g)=|\langle g\rangle|$ divides $n$ by Lagrange and Proposition~\ref{prop:order}(b); then $g^n=(g^{o(g)})^{n/o(g)}=e$. (b) Apply (a) to $(\Z/m\Z)^*$, of order $\varphi(m)$. (c) $o(g)$ divides the prime $n$ and is not $1$, so it is $n$, and $\langle g\rangle$ has $n$ elements: it is $G$.
\end{proof}

So the proof of Euler's theorem given in Volume~1, which multiplied all the invertible residues together, is a special case of a general fact about all finite groups. Lagrange's theorem has no converse: $A_4$, of order $12$, has no subgroup of order $6$ (Chapter~4).

\section{Cyclic groups and isomorphisms}

Two groups can be ``the same'' while their elements have different names. The rotations $\{e,r,r^2,r^3\}$ of a square, the group $\mathbb U_4=\{1,i,-1,-i\}$ and $\Z/4\Z$ have the same multiplication table, up to renaming.

\begin{definition}
An \emph{isomorphism} $f:G\to G'$ is a bijection with $f(ab)=f(a)f(b)$ for all $a,b$. Then $f(e)=e'$, $f(a^{-1})=f(a)^{-1}$, and $f^{-1}$ is an isomorphism too. If one exists, $G$ and $G'$ are \emph{isomorphic}, written $G\cong G'$. A group is \emph{cyclic} if it is generated by one element.
\end{definition}

(Indeed $f(e)=f(ee)=f(e)f(e)$ gives $f(e)=e'$ by cancellation, then $f(a)f(a^{-1})=f(e)=e'$. For $f^{-1}$, apply $f$ to both sides of $f^{-1}(xy)=f^{-1}(x)f^{-1}(y)$.) Isomorphic groups share every property defined from the operation alone: order, being abelian or cyclic, the orders of elements, the number of subgroups of each order.

\begin{theorem}[cyclic groups]\label{thm:cyclic}
Let $G=\langle g\rangle$ be cyclic.
\begin{enumerate}[label=(\alph*)]
\item If $G$ is infinite, $k\mapsto g^k$ is an isomorphism $\Z\to G$. If $|G|=n$, $k\bmod n\mapsto g^k$ is an isomorphism $\Z/n\Z\to G$.
\item Every subgroup of $G$ is cyclic. If $|G|=n$, then for every divisor $d$ of $n$ there is exactly one subgroup of order $d$, namely $\langle g^{n/d}\rangle$, and there are no others.
\item If $|G|=n$, $g^k$ generates $G$ if and only if $\gcd(k,n)=1$. So $G$ has exactly $\varphi(n)$ generators, and exactly $\varphi(d)$ elements of order $d$ for each $d\mid n$.
\end{enumerate}
\end{theorem}
\begin{proof}
(a) The map is surjective by definition of $\langle g\rangle$ (Proposition~\ref{prop:subgroup}(b), with $S=\{g\}$, every product being a power of $g$), and it turns sums into products since $g^{k+l}=g^kg^l$. In the infinite case the powers are distinct (Proposition~\ref{prop:order}), so it is injective. In the finite case it is well defined and injective on classes mod $n$ by Proposition~\ref{prop:order}(a).

(b) Let $H\le G$, $H\ne\{e\}$ (the trivial subgroup is $\langle e\rangle$). Let $m$ be the least $k\ge1$ with $g^k\in H$; it exists since $H$ contains some $g^k\ne e$ and then also $g^{-k}$. For $g^k\in H$, write $k=qm+r$ with $0\le r<m$: then $g^r=g^k(g^m)^{-q}\in H$, so $r=0$ by minimality. Hence $H=\langle g^m\rangle$. If $|G|=n$, then $g^n=e\in H$ gives $m\mid n$, and $|H|=o(g^m)=\frac nm$ by Proposition~\ref{prop:order}(c). So a subgroup of order $d$ must be $\langle g^{n/d}\rangle$, and this one has order $d$.

(c) By Proposition~\ref{prop:order}(c), $o(g^k)=\frac n{\gcd(n,k)}$, which equals $n$ iff $\gcd(n,k)=1$. An element of order $d$ generates a subgroup of order $d$, which is the unique one $C_d$ of (b), cyclic of order $d$; so the elements of order $d$ are the generators of $C_d$, and there are $\varphi(d)$ of them.
\end{proof}

Since every element of a group of order $n$ has an order dividing $n$, counting the elements of $\Z/n\Z$ by their order gives Gauss' identity $\sum_{d\mid n}\varphi(d)=n$, proved in Volume~1 by another route.

\begin{beauty}
Three facts proved separately in Volume~1, Euler's theorem, $\sum_{d\mid n}\varphi(d)=n$ and the counting of primitive roots, are now one-line consequences of two general theorems about groups. This is the economy of abstraction: prove a fact once, in the right generality, and it holds for permutations, matrices, residues and rotations at the same time. Exercise~\ref{ex:fieldcyclic} goes one step further: it proves that $(\Z/p\Z)^*$ is cyclic by an argument that works in every field, and it will give the structure of finite fields in Chapter~9.
\end{beauty}

\section{Exercises}

\begin{exercise}[$\star$]\label{ex:square}
Let $G$ be a group in which $x^2=e$ for every $x$. Show that $G$ is abelian.
\end{exercise}

\begin{exercise}[$\star$]\label{ex:dihedralcentre}
In $D_n$, show that the rotations form a subgroup of index $2$, and find the \emph{centre} $Z(D_n)=\{z:zg=gz\text{ for all }g\}$. (The answer depends on the parity of $n$.)
\end{exercise}

\begin{exercise}[$\star$]\label{ex:order4}
Show that a group of order $4$ is either cyclic or has all its non-identity elements of order $2$. Write the multiplication table in the second case, and give two geometric examples of it.
\end{exercise}

\begin{exercise}[$\star\star$]\label{ex:union}
Show that a group is never the union of two proper subgroups. Give a group that is the union of three proper subgroups.
\end{exercise}

\begin{exercise}[$\star\star$]\label{ex:gl2}
Let $p$ be prime. Show that $|GL_2(\Z/p\Z)|=(p^2-1)(p^2-p)$ and $|SL_2(\Z/p\Z)|=p(p^2-1)$. Check that $GL_2(\Z/2\Z)$ has order $6$ and is not abelian.
\end{exercise}

\begin{exercise}[$\star\star$]\label{ex:cubeS4}
\emph{The cube is $\mathfrak S_4$.} Number the four long diagonals of the cube $1,2,3,4$. Each rotation of the cube permutes the diagonals, which gives a map $\rho$ from the rotation group $C$ of Proposition~\ref{prop:cube} to $\mathfrak S_4$. Show that $\rho$ is an isomorphism.
\end{exercise}

\begin{exercise}[$\star\star$]\label{ex:prodorder}
Let $a,b$ be commuting elements of finite orders $m$ and $n$ with $\gcd(m,n)=1$. Show that $o(ab)=mn$. Show by an example in $\mathfrak S_3$ that the hypothesis $ab=ba$ cannot be dropped.
\end{exercise}

\begin{exercise}[$\star\star\star$]\label{ex:fieldcyclic}
Let $K$ be a field and $G$ a finite subgroup of $K^*$, of order $n$. Show that $G$ is cyclic. \emph{(For $K=\Z/p\Z$ this is the existence of primitive roots, Volume~1, Chapter~5.)}
\end{exercise}

\begin{exercise}[$\star\star\star$]\label{ex:wilson}
\emph{Wilson, from groups.} Let $G$ be a finite abelian group, written multiplicatively.
\begin{enumerate}[label=(\alph*)]
\item Show that the product of all the elements of $G$ equals the product of its elements of order $2$.
\item If $G$ has exactly one element $t$ of order $2$, show that this product is $t$; if it has none, the product is $e$.
\item Deduce Wilson's theorem: $(p-1)!\equiv-1\pmod p$ for every prime $p$.
\end{enumerate}
\end{exercise}

\section{Hints}

\begin{hint}{ex:square}
$x^2=e$ means $x=x^{-1}$. Apply it to $x$, $y$ and $xy$.
\end{hint}
\begin{hint}{ex:dihedralcentre}
Use $sr^k=r^{-k}s$. When does $r^k$ commute with $s$? Can a reflection commute with $r$?
\end{hint}
\begin{hint}{ex:order4}
By Lagrange, the orders of elements are $1$, $2$ or $4$. Without an element of order $4$, use Proposition~\ref{prop:rules}(d) to fill the table. For examples, think of a rectangle, or of $(\Z/8\Z)^*$.
\end{hint}
\begin{hint}{ex:union}
If $G=H\cup K$ with $h\in H\setminus K$ and $k\in K\setminus H$, where is $hk$? For three subgroups, try a group of order $4$.
\end{hint}
\begin{hint}{ex:gl2}
Choose the first column (any non-zero vector), then the second (any vector not on the line of the first). For $SL_2$, the determinant is onto $(\Z/p\Z)^*$, and each value is taken equally often.
\end{hint}
\begin{hint}{ex:cubeS4}
$\rho$ respects products because both sides are composition. Since both groups have $24$ elements, it is enough to show that only the identity fixes every diagonal. Use the signed permutation matrices of the proof of Proposition~\ref{prop:cube}.
\end{hint}
\begin{hint}{ex:prodorder}
$(ab)^k=a^kb^k$. If $(ab)^k=e$, then $a^k=b^{-k}$ lies in $\langle a\rangle\cap\langle b\rangle$. What can the order of an element of this intersection be?
\end{hint}
\begin{hint}{ex:fieldcyclic}
Let $\psi(d)$ be the number of elements of order $d$ in $G$. If $\psi(d)\ne0$, an element of order $d$ generates $d$ distinct roots of $X^d-1$, which has at most $d$ roots in a field. Compare $\psi(d)$ with $\varphi(d)$ and use $\sum_{d\mid n}\varphi(d)=n$.
\end{hint}
\begin{hint}{ex:wilson}
(a) Pair each element with its inverse. (b) The elements of order $2$, together with $e$, form a subgroup. (c) In $(\Z/p\Z)^*$, $x^2=1$ means $p\mid(x-1)(x+1)$.
\end{hint}

\section{Solutions}

\begin{solution}{ex:square}
\idea Commutativity is the statement $xy=yx$, i.e.\ $(xy)^{-1}=\dots$. Inverses reverse the order of a product (Proposition~\ref{prop:rules}(c)), and here every element is its own inverse, so reversing changes nothing.
\writeup For all $x$, $x^2=e$ gives $x^{-1}=x$. So for $x,y\in G$, $xy=(xy)^{-1}=y^{-1}x^{-1}=yx$.
\end{solution}

\begin{solution}{ex:dihedralcentre}
\idea All computations in $D_n$ reduce to the relation $sr^k=r^{-k}s$ (Proposition~\ref{prop:dihedral}). An element commutes with everything iff it commutes with the generators $r$ and $s$. So test the two kinds of elements, $r^k$ and $sr^k$, against $r$ and $s$.
\writeup The rotations $\langle r\rangle$ form a subgroup of order $n$, hence of index $\frac{2n}n=2$ by Lagrange. An element commutes with all of $D_n=\langle r,s\rangle$ iff it commutes with $r$ and $s$.
\begin{itemize}
\item $r^k$ commutes with $r$, and $sr^k=r^{-k}s$, so $r^k$ commutes with $s$ iff $r^k=r^{-k}$, i.e.\ $r^{2k}=e$, i.e.\ $n\mid2k$.
\item $sr^k$ never commutes with $r$: $(sr^k)r=sr^{k+1}$, while $r(sr^k)=sr^{-1}r^k=sr^{k-1}$, and these differ since $r^2\ne e$ for $n\ge3$.
\end{itemize}
So $Z(D_n)=\{e\}$ if $n$ is odd, and $Z(D_n)=\{e,r^{n/2}\}$ if $n$ is even ($r^{n/2}$ is the half-turn, which commutes with every symmetry).
\end{solution}

\begin{solution}{ex:order4}
\idea Lagrange restricts the orders of elements to divisors of $4$. If some element has order $4$ we are done. Otherwise every element is its own inverse, and the ``each element once per row'' rule forces the table.
\writeup By Corollary~\ref{cor:lagrange}, orders are $1,2,4$. If some $g$ has order $4$, $G=\langle g\rangle$ is cyclic. Otherwise $G=\{e,a,b,c\}$ with $a^2=b^2=c^2=e$. Then $ab\notin\{e,a,b\}$: $ab=e$ would give $b=a^{-1}=a$, $ab=a$ would give $b=e$, $ab=b$ would give $a=e$. So $ab=c$, and in the same way the product of any two distinct non-identity elements is the third. The table is
\[
\begin{array}{c|cccc}&e&a&b&c\\\hline e&e&a&b&c\\a&a&e&c&b\\b&b&c&e&a\\c&c&b&a&e\end{array}
\]
This is the \emph{Klein four-group}. Examples: the symmetries of a non-square rectangle (identity, two reflections, the half-turn), and $(\Z/8\Z)^*=\{1,3,5,7\}$, where every square is $1$.
\end{solution}

\begin{solution}{ex:union}
\idea Suppose $G=H\cup K$ with neither contained in the other, and pick witnesses $h\in H\setminus K$, $k\in K\setminus H$. The product $hk$ must land in $H$ or in $K$, and either choice lets us solve for the other witness inside the wrong subgroup.
\writeup Suppose $G=H\cup K$ with $H,K$ proper. If $H\subseteq K$ then $G=K$, a contradiction; so there are $h\in H\setminus K$ and, symmetrically, $k\in K\setminus H$. If $hk\in H$ then $k=h^{-1}(hk)\in H$, a contradiction; if $hk\in K$ then $h=(hk)k^{-1}\in K$, a contradiction. So $hk\notin H\cup K=G$, which is absurd.
The Klein four-group $\{e,a,b,c\}$ of Exercise~\ref{ex:order4} is the union of its three subgroups $\{e,a\}$, $\{e,b\}$, $\{e,c\}$.
\end{solution}

\begin{solution}{ex:gl2}
\idea Counting invertible matrices means counting bases of $(\Z/p\Z)^2$ (the columns), chosen one vector at a time. For $SL_2$, the determinant splits $GL_2$ into $p-1$ classes, and multiplying the first column by a scalar moves between them bijectively.
\writeup A $2\times2$ matrix over the field $\Z/p\Z$ is invertible iff its columns are linearly independent. The first column is any non-zero vector: $p^2-1$ choices. The second is any vector not a multiple of the first: $p^2-p$ choices. So $|GL_2(\Z/p\Z)|=(p^2-1)(p^2-p)$.
For $\lambda\in(\Z/p\Z)^*$, multiplying the first column by $\lambda$ is a bijection from the matrices of determinant $1$ onto those of determinant $\lambda$ (its inverse multiplies by $\lambda^{-1}$). So the $p-1$ values of the determinant are taken equally often, and $|SL_2(\Z/p\Z)|=\frac{(p^2-1)(p^2-p)}{p-1}=p(p^2-1)$.
For $p=2$: $3\cdot2=6$ elements, and $\begin{pmatrix}1&1\\0&1\end{pmatrix}\begin{pmatrix}1&0\\1&1\end{pmatrix}=\begin{pmatrix}0&1\\1&1\end{pmatrix}$ while the product in the other order is $\begin{pmatrix}1&1\\1&0\end{pmatrix}$.
\end{solution}

\begin{solution}{ex:cubeS4}
\idea A map between finite sets of the same size is bijective as soon as it is injective. Injectivity of a product-preserving map comes down to: only the identity goes to the identity permutation (if $\rho(f)=\rho(g)$ then $\rho(g^{-1}f)=\mathrm{id}$). So we must show that a rotation fixing each diagonal (as a line, possibly reversing it) is the identity. The matrix description makes this a small computation.
\writeup For rotations $f,g$, the permutation $\rho(fg)$ of the diagonals is ``apply $g$, then $f$'', which is $\rho(f)\rho(g)$. If $\rho(f)=\rho(g)$ then $\rho(g^{-1}f)=\rho(g)^{-1}\rho(f)=\mathrm{id}$, so it suffices to show that a rotation $h$ with $\rho(h)=\mathrm{id}$ is the identity.
Such an $h$ maps each diagonal to itself, so $h(v)=\pm v$ for each of the vertices $v=(1,1,1)$, $(1,1,-1)$, $(1,-1,1)$. By the proof of Proposition~\ref{prop:cube}, $h$ is a signed permutation matrix with determinant $1$. The coordinates of $h(1,1,1)$ are the non-zero entries of the rows of $h$; since $h(1,1,1)=\varepsilon(1,1,1)$, all these entries equal $\varepsilon$, and $h=\varepsilon P$ for a permutation matrix $P$ (sending $e_j$ to $e_{\sigma(j)}$). Then $P(1,1,-1)$ has its entry $-1$ in position $\sigma(3)$, and it must be $\pm(1,1,-1)$; since it has a single $-1$, it is $(1,1,-1)$ and $\sigma(3)=3$. In the same way $(1,-1,1)$ gives $\sigma(2)=2$, hence $\sigma=\mathrm{id}$ and $h=\varepsilon I_3$. Since $\det h=\varepsilon^3=1$, $h=I_3$.
So $\rho$ is injective, and as $|C|=|\mathfrak S_4|=24$, it is bijective: an isomorphism.
\end{solution}

\begin{solution}{ex:prodorder}
\idea Since $a,b$ commute, $(ab)^k=a^kb^k$, so $(ab)^{mn}=e$ and $o(ab)\mid mn$. For the converse, $(ab)^k=e$ makes $a^k=b^{-k}$ an element of both cyclic subgroups, whose orders are coprime, so it must be $e$.
\writeup As $ab=ba$, $(ab)^k=a^kb^k$ for all $k$ (induction). Suppose $(ab)^k=e$. Then $x=a^k=b^{-k}$ lies in $\langle a\rangle\cap\langle b\rangle$, so $o(x)$ divides $m$ and $n$ (Lagrange in both cyclic groups), hence $o(x)=1$. So $a^k=e$ and $b^k=e$, i.e.\ $m\mid k$ and $n\mid k$, hence $mn\mid k$ since $\gcd(m,n)=1$. Conversely $(ab)^{mn}=(a^m)^n(b^n)^m=e$. So $o(ab)=mn$.
In $\mathfrak S_3$, $a=(1\,2)$ has order $2$ and $b=(1\,2\,3)$ has order $3$, but $ab$ sends $1\mapsto 2\mapsto 1$, $2\mapsto3\mapsto3$, $3\mapsto1\mapsto2$: $ab=(2\,3)$ has order $2$, not $6$.
\end{solution}

\begin{solution}{ex:fieldcyclic}
\idea We want an element of order $n$. Count elements by their order: $\psi(d)$ elements of order $d$, with $\sum_{d\mid n}\psi(d)=n$ (Lagrange). A field enters only through ``a polynomial of degree $d$ has at most $d$ roots'' (Volume~0, Chapter~3), which caps $\psi(d)$ by $\varphi(d)$. Since the $\varphi(d)$ also sum to $n$, every cap is reached, in particular $\psi(n)=\varphi(n)\ge1$.
\writeup For $d\mid n$, let $\psi(d)$ be the number of elements of order $d$ in $G$. By Corollary~\ref{cor:lagrange}, every element has order dividing $n$, so $\sum_{d\mid n}\psi(d)=n$.
Suppose $\psi(d)\ge1$, and let $x$ have order $d$. The $d$ distinct elements $1,x,\dots,x^{d-1}$ are roots of $X^d-1$, which has at most $d$ roots in the field $K$. So every element of order $d$, being a root of $X^d-1$, lies in $\langle x\rangle$, a cyclic group of order $d$, which has exactly $\varphi(d)$ elements of order $d$ (Theorem~\ref{thm:cyclic}(c)). Hence $\psi(d)\in\{0,\varphi(d)\}$ for every $d$, so $\psi(d)\le\varphi(d)$. Now
\[
n=\sum_{d\mid n}\psi(d)\le\sum_{d\mid n}\varphi(d)=n,
\]
so equality holds term by term. In particular $\psi(n)=\varphi(n)\ge1$: $G$ has an element of order $n$ and is cyclic.
\end{solution}

\begin{solution}{ex:wilson}
\idea In the product of all elements, each element can be paired with its inverse, and each pair contributes $e$ (commutativity lets us reorder). Only the elements equal to their own inverse are unpaired: $e$ and the elements of order $2$. Then for (b), the set of these is a subgroup, and the product over a subgroup is easier to see by pairing once more.
\writeup (a) Since $G$ is abelian the product $P=\prod_{g\in G}g$ does not depend on the order of the factors. The map $g\mapsto g^{-1}$ is a bijection of $G$ with itself; it pairs each $g$ with $g\ne g^{-1}$ to its inverse, and these pairs contribute $e$. The fixed points are the $g$ with $g^2=e$: $e$ and the elements of order $2$. So $P$ is the product of the elements of order $2$.
(b) If $t$ is the only element of order $2$, then $P=t$. If there is none, $P=e$. (More generally, the set $T=\{g:g^2=e\}$ is a subgroup, since $(gh)^2=g^2h^2=e$.)
(c) In $G=(\Z/p\Z)^*$ (abelian, of order $p-1$), $x^2=1$ means $p\mid(x-1)(x+1)$, so $x=\pm1$ by Euclid's lemma. For $p$ odd, $-1$ is the only element of order $2$, so the product of all classes $1,2,\dots,p-1$ is $-1$: $(p-1)!\equiv-1\pmod p$. For $p=2$, $1!=1\equiv-1\pmod 2$.
\end{solution}

\section*{Sources and further reading}

\begin{itemize}
\item M.~Artin, \emph{Algebra}, 2nd ed., Pearson, 2011, Chapters~2 and~6. Groups introduced through symmetry, with the dihedral group and the rotation groups of the Platonic solids.
\item M.~A.~Armstrong, \emph{Groups and Symmetry}, Springer, 1988. A gentle and geometric first course, from which the cube and its long diagonals are borrowed.
\item N.~Carter, \emph{Visual Group Theory}, Mathematical Association of America, 2009. Cayley diagrams and multiplication tables as pictures, in the spirit of Figure~\ref{fig:symmetries}.
\item D.~S.~Dummit and R.~M.~Foote, \emph{Abstract Algebra}, 3rd ed., Wiley, 2004, Chapters~1--3. The standard reference for everything in this chapter, with many more exercises.
\item W.~Burnside, \emph{Theory of Groups of Finite Order}, 2nd ed., Cambridge University Press, 1911 (Dover reprint, 1955). The first great book on finite groups, still pleasant to browse.
\item T.~Rokicki, H.~Kociemba, M.~Davidson and J.~Dethridge, \emph{The diameter of the Rubik's cube group is twenty}, SIAM J.\ Discrete Math.\ 27 (2013), 1082--1105 (the computation was announced in 2010).
\end{itemize}
